\scanpage{10}
\unsure{deuxième|\deleted{deuxième}} élément de structure
\[
f:\widetilde{S}\longrightarrow\unsure{\mathscr{J}|\mathcal{J}|\gamma}
\]
Ces deux éléments de structure sont reliés par la condition que $f$ est compatible avec les actions de $\pi$, quand on le fait agir sur $\mathscr{J}$ via $\alpha:\pi\dashvG(\mathbb{Q})\subset G(\mathbb{R})$ et l'action \inserted{à \uline{droite}} de $G(\mathbb{R})$ sur $\mathscr{J}$ (opposée à l'action gauche).

La donnée d'un couple $(F,\varphi)$ équivaut à celle d'un couple $(f,\alpha)$ satisfaisant les deux conditions précédentes.

\deleted{Il faudrait fixer des \uline{réseaux}, et des sous-groupes discrets convenables de $G(\mathbb{Q})$, du type $G(\mathbb{Z})$ commensurable à un tel ....}

Autre façon de présenter les ``choses, \unsure{indépend.ment|indépendamment} du choix d'un point-base $s$ de $S$. \deleted{(\unsure{reprenons|recommençons})}.

\deleted{On dit que $F$ est co} Notons que les foncteurs fibres \inserted{$\Phi_s$ (rel/$\mathbb{Q}$)} sur \deleted{str. de Hodge} $\mathrm{Hodge}(S)$ définis par les divers points $s\in S$ \deleted{\unsure{sont}} \inserted{d'une \unsure{$\overline{m}$|\deleted{c}} composante connexe $S^{i}$ de $S$ sont} isomorphes ; donc aussi les $\Phi_s\circ F$ \inserted{$(s\in S^{i})$}. On dit que \inserted{la $G$-subordination} $F$ est \uline{compatible avec le foncteur} \uline{oubli} \uline{sur} $C=\operatorname{Rep}(G)$ si les $\Phi_s\circ F$ sont isomorphes audit foncteur oubli. Soit alors, pour
% archive sheet number: 9

